% Fichier complété par tools/complete_missing_corrections.py.
% Source : STAT/STAT-02-Local/39_SeineMusicale/39_SeineMusicale.tex
% Statut : rédigé (4 question(s)).
\begin{corrigebox}[Corrigé rédigé]
\CorrigeQuestion{1}
La pression dynamique est $p=f$ (notation de l'énoncé). Sur
                $dS$, $d\vec F_{vent}=-f\,dS\,\vec x_{C_G}$, normale à la voile et
                opposée au vent.
\CorrigeQuestion{2}
On paramètre la demi-voile par des coordonnées polaires. Le moment
                élémentaire est $dM_z=(\overrightarrow{OP}\wedge d\vec F)\cdot\vec z$;
                l'intégration sur le demi-disque conduit à un terme en $fR^3$ multiplié
                par la fonction trigonométrique de $\alpha$ issue de la projection de la
                normale.
\CorrigeQuestion{3}
Par définition,
                $F_{vent}=M_{O,z}/[(\overrightarrow{OC_G}\wedge\vec x_{C_G})\cdot\vec z]$.
                On remplace $M_{O,z}$ par le résultat précédent et
                $OC_G=4R/(3\pi)$ pour un demi-disque homogène.
\CorrigeQuestion{4}
Le maximum est obtenu en annulant la dérivée de la fonction
                trigonométrique en $\alpha$ et en comparant les bornes. Pour une dépendance
                en $\sin\alpha\cos\alpha$, il est atteint à $\alpha=\pi/4$ et vaut la
                moitié du coefficient multiplicatif.
\end{corrigebox}
